% Generated draft: requires mathematical and source review.
\documentclass[11pt]{article}
\usepackage[margin=25mm]{geometry}
\usepackage{fontspec}
\setmainfont{Latin Modern Roman}
\usepackage{amsmath,amssymb,mathtools,unicode-math}
\setmathfont{Latin Modern Math}
\usepackage{enumitem,needspace,xurl,hyperref}
\setlength{\parindent}{0pt}
\setlength{\parskip}{5pt}
\setlength{\emergencystretch}{3em}
\newcommand{\N}{\mathbb{N}}
\newcommand{\Z}{\mathbb{Z}}
\newcommand{\Q}{\mathbb{Q}}
\newcommand{\R}{\mathbb{R}}
\newcommand{\C}{\mathbb{C}}
\newcommand{\F}{\mathbb{F}}
\newcommand{\A}{\mathbb{A}}
\newcommand{\PP}{\mathbb{P}}
\newcommand{\E}{\mathbb{E}}
\newcommand{\eps}{\varepsilon}
\newcommand{\dd}{\,\mathrm{d}}
\begin{document}
\section*{UnsolvedMath ID: 4000016}
\textbf{Source problem code:} catalogue entry
\subsection{Definitions and mathematical statement}
Let $(X,d)$ be a metric space with probability transition laws $m_x$, and let $\Pi(\mu,\nu)$ denote the set of couplings of two probability measures. Define the extended $L^\infty$ transport distance by
\[
 T_\infty(\mu,\nu)=\inf_{\pi\in\Pi(\mu,\nu)}
          \operatorname*{ess\,sup}_{(u,v)\sim\pi} d(u,v).
\]
The source proposes replacing the $T_1$ transport in coarse Ricci curvature by this worst-displacement notion. Its explicitly described zero-threshold comparison requires, for distinct $x,y$, a coupling $\pi\in\Pi(m_x,m_y)$ satisfying
\[
 \pi\bigl(\{(u,v):d(u,v)\le d(x,y)\}\bigr)=1.
\]
This spells out the measure-theoretic meaning of moving every point by at most $d(x,y)$. Without assumptions ensuring attainment of the infimum, the existence of such a coupling should not simply be identified with the numerical inequality $T_\infty(m_x,m_y)\le d(x,y)$.

A direct numerical analogue of the one-step coarse-Ricci formula would be
\[
 \kappa_\infty(x,y)=1-\frac{T_\infty(m_x,m_y)}{d(x,y)}.
\]
This is a candidate normalization illustrating the proposed replacement; the source asks whether a numerical value can be assigned and does not specify this formula as its final definition. Infinitesimal scales, alternative normalizations, and any treatment of an infinite transport distance are not fixed.

Ollivier's Problem P asks whether this $L^\infty$ coupling approach gives a useful notion of discrete sectional curvature, how to quantify it, and whether it is related to Alexandrov sectional curvature. A precise Alexandrov comparison would have to specify the metric-space class, the geodesic or triangle-comparison hypotheses, lower versus upper curvature bounds, the choice and scale of transition laws, and the direction or limiting meaning of the proposed relation. None of these comparison data is supplied, so no particular equivalence or implication is asserted here.
\subsection{Short English statement}
Study a discrete sectional-curvature notion based on controlling the largest displacement in a coupling of neighboring transition laws, rather than their average transport distance. Determine its numerical formulation and its relationship to Alexandrov curvature.

The supplied summary reports functional-analytic characterizations and applications of an $L^\infty$-based notion, while describing the exact Alexandrov connection as unresolved. It provides no precise result or later citation for those developments. The partial-progress label and original 2008 open formulation are supplied assessments, with no later authoritative resolution check recorded.
\subsection{Sources}
\begin{itemize}
\item Local UnsolvedMath registry, permanent ID 4000016, source code catalogue entry, ``Discrete sectional curvature,'' proposed by Yann Ollivier, 2008. Dataset curation and metadata: Ulam AI and UnsolvedMath Contributors, CC BY 4.0.
\item Yann Ollivier, \emph{Discrete Ricci curvature: Open problems} (2008), Problem P, \url{https://www.yann-ollivier.org/rech/publs/problems_curvmarkov.pdf}; pdf-text extraction, accessed 2026-07-29.
\item The supplied research summary mentions functional-analytic characterizations and applications without identifying primary papers, hypotheses, or theorem statements. No additional reference is invented.
\item Source provenance marks status evidence NEEDS\_REVIEW and rights NEEDS\_REVIEW: the source presents the item as open, no later authoritative resolution check is recorded, and redistribution permission was not established despite verified public access. This qualification is separate from the dataset licence.
\end{itemize}
\end{document}
